\begin{exercisechunk}
\item \textbf{When the AGOP recovers a multi-index subspace.}
\label[exercise]{ex:l15-agop-subspace}
\exerciseprofile{2}{2}{2}{\exproof\quad\excalculation\quad\exconstruction\quad\exinterpretation}{%
Prove the population identity behind gradient-based dimension reduction and
identify the nondegeneracy condition required for subspace recovery.}
Let $X$ be a random vector in $\R^d$, let $W\in\R^{d\times s}$ have full
column rank, and suppose $f(x)=g(W^\top x)$ with differentiable $g$. Assume
the displayed expectations below are finite.
\begin{enumerate}[(a),beginpenalty=10000]
\item Use the chain rule to prove
\[
 G_f:=\E[\nabla f(X)\nabla f(X)^\top]
 =W G_g W^\top,
 \qquad
 G_g:=\E[\nabla g(W^\top X)\nabla g(W^\top X)^\top].
\]
\item Prove that $\operatorname{range}(G_f)\subseteq
\operatorname{range}(W)$. Show that equality holds
when $G_g$ is positive definite.
\item Give an example in which $W$ has two columns but $G_f$ has rank one.
What does the example say about recovering a declared index space from
gradients? Explain why an arbitrary interpolant can have a different AGOP
from the regression target and why a recursive feature machine therefore
needs the AGOP of an accurate predictor.
\end{enumerate}
\end{exercisechunk}
